Work first over $R_0=\mathbb Z[1/210]$. The effective level-change group
has order $36$, prime to every residue characteristic of $R_0$. Its generic
action is free. For a finite tame group acting on a smooth relative curve,
the fixed locus of an element is smooth locally along its fixed points;
an isolated special fixed point consequently lifts. Thus no new fixed
point occurs here, and $Y_{R_0}\to X_{R_0}$ is finite etale. The $W$ action
extends to the smooth stable model of $X$. Its stabilizers remain cyclic,
of order at most two. Distinct involution fixed loci cannot intersect;
each fixed locus is etale over the base. Consequently $X_{R_0}/W$ is a
smooth genus-zero curve with five disjoint branch sections of index two.
The three rational sections already chosen identify it with
$\mathbb P^1_{R_0}$. Their values are the five elements
of $\mathcal C$.

Set $K=\widehat{\mathcal O}^{\times}$, $U=U(3)$,
$Z=\mathbb A_f^\times$, and $N=N_{B(\mathbb A_f)^\times}
(\widehat{\mathcal O})$. Normalization preserves $3\widehat{\mathcal O}$,
so $U$ is normal in $N$. The canonical adelic action is defined over the
reflex field $\mathbb Q$ \cite[2.2.5 and 2.3.1]{Deligne1979}; this refers to the entire canonical model, not
an individually selected complex component. Every central finite idele
is a rational scalar times a unit idele. At three its scalar residue is
$1$ or $-1$, which a rational sign removes. Thus $Z\subset\mathbb Q^\times U$
acts trivially on $Y$.

At three the quaternion algebra is ramified. Its maximal order modulo
three has $81$ elements, radical of size $9$, and hence $72$ units.
It follows that
\[
 1\longrightarrow (\mathcal O/3\mathcal O)^\times/\{\pm1\}
 \longrightarrow\Gamma=N/(ZU)\longrightarrow W\longrightarrow1,
 \qquad |\Gamma|=36\cdot16=576.
\]
Above an unbranched point the actions of $W$ and of the effective
level-change kernel are free. The fiber therefore is a torsor of size
$576$ for this constant group. The Galois permutations commute with its
regular action, so belong to its regular centralizer, also of order
$576$. The normal splitting field $L$ consequently satisfies
$[L:\mathbb Q]\mid576$: the Galois group is a subgroup of that
regular centralizer. In particular $[L:\mathbb Q]\le576$.

All fiber surfaces and their three-torsion markings are defined over $L$.
For two such surfaces, a homomorphism annihilating the three-torsion
factors through $[3]$. Hence reduction embeds
\[
 \operatorname{Hom}(A,B)/3\operatorname{Hom}(A,B)
 \hookrightarrow\operatorname{Hom}(A[3],B[3]).
\]
The finite Galois image on the integral Hom lattice is trivial modulo
three. The principal congruence subgroup modulo three in an integral
general linear group is torsion-free, so this image is trivial. Thus
every geometric homomorphism between fiber surfaces is defined over $L$.

At a prime outside $210$, tame local coordinates for the quotient are
$w=z^2$ up to a unit at a branch section, and etale elsewhere. A base
point with even contact gives no ramification after an unramified
extension. The universal abelian scheme extends at the resulting
integral point of $Y_{R_0}$. Therefore
\[
 \operatorname{rad}(D_LN_{\rm bad})\mid B_*M(s).
\tag{R}\label{eq:h5-R}
\]
This assertion includes contact with infinity. It is not a statement
obtained merely from the numerical prime support of pairwise determinants.

The normalizer correspondences give isogenies between the conjugate
surfaces. In the non-CM case the geometric endomorphism algebra is $B$;
Skolem--Noether adjusts a conjugate isogeny to commute with its chosen
quadratic subfield. All such maps lie over $L$ by the Hom argument.
These facts establish the usual GL$_2$-type and compatibility hypotheses;
they do \emph{not} establish the separate bounded descent to a
GL$_2$-type variety over $\mathbb Q$.
